\scanpage{146}
Ainsi, le groupe quotient de $D_{\mathbb R\times\mathbb R}$ qui opère fidèlement sur $X$ est le groupe
\[
(21)\qquad D_{\mathbb R\times\mathbb U}\overset{\mathrm{dfn}}{=}\{1,\sigma\}.(\mathbb R\times\mathbb U)\qquad(1/2\text{ direct})
\]
où $\sigma$ opère \inserted{encore} sur $\mathbb R\times\mathbb U$ par symétrie
\[
\sigma(u,\theta)=(-u,\theta^{-1}).
\]
On fait opérer \deleted{\unsure{ce}} \inserted{$D_{\mathbb R\times\mathbb U}$} sur $X$ par
\[
(22)\qquad\left\{\begin{aligned}
\sigma(z,s)&=(z^{-1},s)\\
g_{\underset{\substack{\rotatebox{90}{$\in$}\\\mathbb R}}{a},\underset{\substack{\rotatebox{90}{$\in$}\\\mathbb U}}{\theta}}(z,s)&\overset{\text{déf}}{=}(z\exp(2i\pi as)\theta,s)
\end{aligned}\right.
\]
et cette action se relève en une action de $D_{\mathbb R\times\mathbb R}$ sur $\widetilde X$ par
\[
(23)\qquad\left\{\begin{aligned}
\sigma(t,s)&\overset{\text{déf}}{=}(-t,s)\\
f_{\underset{\substack{\rotatebox{90}{$\in$}\\\mathbb R}}{a},\underset{\substack{\rotatebox{90}{$\in$}\\\mathbb R}}{b}}(t,s)&\overset{\text{déf}}{=}(t+(as+b),s)
\end{aligned}\right.
\]
\deleted{Considérons} On fera opérer le groupe
\[
(24)\qquad\{\pm1\}\unsure{\ltimes|\times}D_{\mathbb R\times\mathbb U}
\]
\margin{(1/2 direct, via op. de $\pm1$ sur $D_{\mathbb R\times\mathbb U}$ par chgt de signe sur le facteur $\mathbb R$}
sur $X$, par
\[
(2\unsure{5|4})\qquad (\varepsilon\times\sigma^\nu.g_{a,\theta})(z,s)\overset{\text{déf}}{=}\left((z^{\deleted{\unsure{-1}}}\exp(2i\pi as)\theta)^{(-1)^\nu},\varepsilon s\right)
\]
cette opération se relève en une op. de
\[
(26)\qquad\{\pm1\}\unsure{\ltimes|\times}D_{\mathbb R\times\mathbb R}
\]
\margin{(1/2 direct, via op. de $\pm1$ sur $D_{\mathbb R\times\mathbb R}$ par chgt de signe sur le facteur $a$ de $f_{a,b}$ ...)}
sur $\widetilde X$, par
\[
(27)\qquad (\varepsilon\times\sigma^\nu f_{\underset{\substack{\rotatebox{90}{$\in$}\\\mathbb R}}{a},\underset{\substack{\rotatebox{90}{$\in$}\\\mathbb R}}{b}})(t,s)\overset{\text{déf}}{=}\left((\unsure{-1|+1})^\nu(t+(as+b)),\unsure{\varepsilon}s\right)
\]
\margin{146}
